\section{Introducción}

\subsection{Planteamiento del Problema}

En los establecimientos del sector restaurantero, la gesti\'on de pedidos y el proceso
de pago representan una parte fundamental de la experiencia del cliente y del control
operativo del negocio. Sin embargo, en muchos restaurantes estos procesos a\'un dependen
de m\'etodos tradicionales que implican la intervenci\'on constante del personal, registros
manuales o sistemas parcialmente digitalizados \cite{toast_success_2023, nra_state_2023}. Esto genera inconvenientes como errores
en los pedidos, retrasos en la comunicaci\'on con cocina y dificultades al momento de dividir
la cuenta entre varios comensales.

Uno de los problemas m\'as frecuentes ocurre cuando un grupo de personas comparte una
mesa y desea pagar \'unicamente por los productos que consumi\'o. En estos casos, el personal
del restaurante debe identificar manualmente los consumos individuales, calcular proporciones
o dividir la cuenta total, lo cual resulta tedioso, propenso a errores y genera tiempos de espera
innecesarios \cite{lee_digital_2024, tan_tabletop_2020}.

En M\'exico, el sector de servicios de preparaci\'on de alimentos y bebidas concentra m\'as de
700{,}000 unidades econ\'omicas de acuerdo con el Directorio Estad\'istico Nacional de Unidades
Econ\'omicas (DENUE) del Instituto Nacional de Estad\'istica y Geograf\'ia (INEGI) \cite{inegi_denue_2024}.
Asimismo, la C\'amara Nacional de la Industria de Restaurantes y Alimentos Condimentados (CANIRAC)
proyect\'o un crecimiento del 4\% para 2024 \cite{inegi_canirac_2023} y de entre 5.5\% y 6\% para 2025 \cite{ramirez_canirac_2025},
impulsado por la recuperaci\'on del consumo y la apertura de nuevas unidades.
A pesar de esta expansi\'on, la brecha tecnol\'ogica entre los sistemas existentes y las expectativas
de los comensales sigue siendo significativa: pocas soluciones integran en una misma plataforma
la interacci\'on directa del comensal con el men\'u, la gesti\'on de \'ordenes y la divisi\'on
aut\'onoma de pagos.

En los \'ultimos a\~nos han surgido diferentes sistemas tecnol\'ogicos orientados a la gesti\'on
de restaurantes, tales como sistemas de punto de venta (POS), aplicaciones de pedidos digitales
o plataformas de administraci\'on de inventarios. No obstante, muchas de estas soluciones se
centran principalmente en la operaci\'on interna del restaurante y no integran completamente
la participaci\'on directa del cliente en el proceso de selecci\'on de productos, gesti\'on de
pedidos y divisi\'on de pagos dentro de una misma plataforma \cite{katagri_smart_2025} de manera individual o colaborativa.

Para efectos del presente trabajo, se entiende la omnicanalidad como la integración
coordinada de las interfaces utilizadas por los distintos actores del restaurante:
comensal, mesero, cocina y administración. Estas interfaces no operan como sistemas
aislados, sino como puntos de interacción conectados a un mismo flujo transaccional,
donde cada acción genera eventos persistidos en una base de datos común y posteriormente
aprovechados para análisis operativo y generación de métricas de negocio.

\subsection{Objetivos}

\subsubsection{Objetivo General}

Diseñar y desarrollar una plataforma web progresiva (PWA) para el sector restaurantero,
orientada a la gestión estructurada de mesas, órdenes y liquidaciones internas, con el
propósito de generar datos transaccionales consistentes que alimenten un modelo de base
de datos centralizado y un pipeline analítico batch basado en Apache Spark.

El sistema permitirá que comensales, meseros, cocina y personal administrativo interactúen
dentro de un flujo operativo coordinado mediante mesas virtuales accesibles por código QR,
registro de órdenes individuales y compartidas, y control interno del estado de liquidación
de cuentas. A partir de estos datos, se implementará una arquitectura analítica tipo
Medallion (Bronze, Silver y Gold) para la transformación, agregación y visualización de
métricas de negocio, tales como ventas por sucursal, productos más consumidos, tiempos
operativos, comportamiento de consumo y estimaciones exploratorias de demanda.

El núcleo técnico del proyecto se centra en el diseño de la base de datos transaccional,
la trazabilidad de los eventos generados durante la operación restaurantera, el procesamiento
analítico mediante Apache Spark y la construcción de módulos de inteligencia de negocio
que apoyen la toma de decisiones administrativas.

\newpage

\subsubsection{Objetivos Específicos}

Los siguientes objetivos específicos describen los módulos funcionales, técnicos y
analíticos que deberán desarrollarse para alcanzar el objetivo general. La plataforma
operativa de mesas, órdenes y liquidaciones internas funcionará como mecanismo de
captura estructurada de datos transaccionales, mientras que el núcleo técnico del
trabajo se centrará en el diseño de la base de datos, la trazabilidad de eventos, el
procesamiento batch con Apache Spark y la generación de indicadores analíticos para
la toma de decisiones.

\noindent\textbf{Módulo 1 — Infraestructura, Base de Datos y Arquitectura Analítica}

\noindent\textbf{OE-01 Infraestructura Técnica del Sistema.}
Configurar la infraestructura base del proyecto utilizando Vercel para el despliegue
de la aplicación web progresiva, Supabase como plataforma de persistencia transaccional
basada en PostgreSQL y un servidor privado virtual (VPS) para la ejecución del pipeline
analítico batch con Apache Spark. Asimismo, establecer los ambientes de desarrollo,
preview y producción necesarios para validar el funcionamiento del sistema.

\textit{Evidencia: Aplicación desplegada en Vercel. Instancia de Supabase activa.
VPS configurado para ejecución de jobs analíticos. Variables de entorno y ambientes
separados documentados.}

\noindent\textbf{OE-02 Diseño de la Base de Datos Transaccional.}
Diseñar e implementar un modelo de base de datos relacional que permita representar
la operación restaurantera mediante entidades como cadena, zona, sucursal, usuarios,
mesas físicas, mesas virtuales, comensales, productos, órdenes, estados de cocina y
registros de liquidación interna. El modelo deberá garantizar integridad referencial,
consistencia transaccional y trazabilidad histórica de los eventos relevantes.

\textit{Evidencia: Schema relacional implementado mediante Prisma ORM. Migraciones
ejecutadas correctamente. Relaciones entre entidades validadas. Restricciones de
integridad documentadas.}

\noindent\textbf{OE-03 Arquitectura Analítica Medallion.}
Definir e implementar una arquitectura analítica basada en capas Bronze, Silver y Gold
para transformar los datos operativos del sistema en conjuntos analíticos. La capa
Bronze almacenará extracciones crudas de la operación; la capa Silver contendrá datos
limpios, normalizados y enriquecidos; y la capa Gold concentrará agregados listos para
ser consultados desde el dashboard administrativo.

\textit{Evidencia: Estructura de almacenamiento Medallion creada. Datasets Bronze,
Silver y Gold generados. Documentación del flujo de transformación entre capas.}

\noindent\textbf{Módulo 2 — Organización, Usuarios y Control de Acceso}

\noindent\textbf{OE-04 Estructura Organizacional del Restaurante.}
Implementar los módulos de creación, edición y consulta de la estructura organizacional
del sistema, considerando tres niveles principales: cadena restaurantera, zona y sucursal.
Cada entidad deberá estar asociada jerárquicamente a su nivel superior, permitiendo que
los datos operativos y analíticos puedan consultarse según el alcance correspondiente.

\textit{Evidencia: CRUD funcional de cadena, zonas y sucursales. Validación de
integridad jerárquica. Consultas filtradas por nivel organizacional.}

\noindent\textbf{OE-05 Autenticación, Roles y Alcance de Usuarios.}
Implementar un sistema de autenticación y autorización para usuarios colaboradores,
definiendo roles con alcances diferenciados: Admin Cadena, Gerente Zona, Gerente
Sucursal, Mesero y Cocina/Chef. El acceso a los datos deberá restringirse según el rol
y el nivel organizacional asignado, aplicando el principio de mínimo privilegio.

\textit{Evidencia: Matriz de roles y permisos documentada. Flujo de inicio de sesión
funcional. Validación de acceso restringido por rol y nivel organizacional.}

\noindent\textbf{OE-06 Gobernanza y Seguridad de Datos.}
Configurar políticas de seguridad a nivel de fila mediante Row Level Security (RLS)
en Supabase para evitar accesos indebidos entre cadenas, zonas y sucursales. Además,
implementar mecanismos de auditoría para registrar operaciones críticas realizadas
sobre entidades relevantes del sistema.

\textit{Evidencia: Políticas RLS activas. Pruebas de aislamiento de datos entre niveles
organizacionales. Tabla de auditoría con registros de usuario actor, entidad afectada,
acción realizada y fecha de operación.}

\noindent\textbf{Módulo 3 — Catálogo y Datos Operativos}

\noindent\textbf{OE-07 Gestión de Catálogo de Productos.}
Desarrollar el módulo de catálogo para administrar categorías, productos, variantes de
precio y modificadores. El catálogo deberá alimentar directamente el flujo de órdenes,
permitiendo capturar datos consistentes sobre productos consumidos, precios aplicados
y configuraciones seleccionadas por los comensales.

\textit{Evidencia: CRUD funcional de categorías, productos, variantes y modificadores.
Productos disponibles para selección desde el menú digital. Datos de catálogo vinculados
correctamente con órdenes registradas.}

\noindent\textbf{OE-08 Captura Histórica de Productos en Órdenes.}
Implementar un mecanismo de snapshots inmutables para los productos registrados dentro
de una orden, almacenando información como nombre, precio, variante y modificadores al
momento de la compra. Esto permitirá conservar la integridad histórica de las órdenes
aun cuando el catálogo cambie posteriormente.

\textit{Evidencia: Orden generada con snapshot de producto. Modificación posterior del
catálogo sin afectar órdenes históricas. Validación de consistencia entre orden,
producto original y datos almacenados.}

\noindent\textbf{Módulo 4 — Mesas Físicas, Mesas Virtuales y Comensales}

\noindent\textbf{OE-09 Gestión de Mesas Físicas.}
Implementar el módulo de administración de mesas físicas por sucursal, permitiendo
registrar número, capacidad, ubicación y estado operativo de cada mesa. El sistema
deberá permitir asociar una o más mesas físicas a una sesión operativa para representar
grupos de comensales durante el servicio.

\textit{Evidencia: Registro y edición de mesas físicas. Visualización de estado de mesa.
Asociación de mesas físicas a sesiones operativas. Persistencia correcta de cambios
en base de datos.}

\noindent\textbf{OE-10 Mesa Virtual y Acceso mediante QR.}
Desarrollar el flujo de inicialización de mesa virtual mediante códigos QR generados
por el personal autorizado. El enlace del QR deberá incluir una firma criptográfica
basada en HMAC-SHA256 para evitar manipulación del identificador de mesa. Al escanear
el QR, el comensal podrá unirse a una mesa virtual activa mediante una sesión temporal,
siempre que la mesa se encuentre disponible y la validación de acceso sea correcta.

\textit{Evidencia: QR generado con firma HMAC-SHA256 válida. Rechazo de enlaces alterados.
Creación de sesión temporal de comensal. Unión exitosa a mesa virtual activa. Rechazo
de acceso cuando la mesa se encuentre inactiva, liquidada o cerrada.}

\noindent\textbf{OE-11 Estados de Mesa Virtual.}
Definir e implementar el ciclo de vida de una mesa virtual mediante estados operativos
claros: \textit{Libre}, \textit{Activa}, \textit{En consumo}, \textit{Por liquidar},
\textit{Liquidada} y \textit{Cerrada}. Los cambios de estado deberán registrarse en
la base de datos para permitir trazabilidad operativa y posterior análisis de tiempos
de atención.

\textit{Evidencia: Transiciones de estado implementadas. Registro de cambios de estado.
Validación de reglas, como impedir nuevas órdenes en mesas liquidadas o cerradas.}

\noindent\textbf{Módulo 5 — Órdenes, Cocina y Eventos Transaccionales}

\noindent\textbf{OE-12 Registro de Órdenes Individuales y Compartidas.}
Implementar la interfaz de menú digital para que los comensales puedan seleccionar
productos, configurar modificadores y registrar órdenes individuales o compartidas.
Cada orden deberá asociarse a la mesa virtual, al comensal correspondiente y al tipo
de consumo, generando datos transaccionales estructurados para su posterior análisis.

\textit{Evidencia: Orden individual registrada correctamente. Orden compartida asociada
a mesa grupal. Relación entre comensal, mesa, productos y tipo de orden validada en
base de datos.}

\noindent\textbf{OE-13 Tablero de Cocina y Estados de Preparación.}
Desarrollar la vista de cocina para visualizar órdenes activas de la sucursal, agrupadas
por estado de preparación. El personal de cocina podrá actualizar el estado de cada
orden, generando eventos operativos que permitan analizar tiempos de preparación,
flujo de trabajo y desempeño de estaciones.

\textit{Evidencia: Nueva orden visible en tablero de cocina. Cambio de estado reflejado
en tiempo real. Registro histórico de transiciones de preparación. Datos disponibles
para análisis posterior.}

\noindent\textbf{OE-14 Sincronización en Tiempo Real.}
Integrar mecanismos de actualización en tiempo real para reflejar cambios de órdenes,
estados de cocina y estados de mesa entre las interfaces del comensal, mesero, cocina
y administración. Esta sincronización permitirá mantener consistencia operativa durante
la sesión activa.

\textit{Evidencia: Actualización de órdenes sin recargar la página. Cambio de estado
de cocina visible para los usuarios correspondientes. Validación de sincronización
entre al menos dos interfaces conectadas simultáneamente.}

\noindent\textbf{Módulo 6 — Liquidación Interna}

\noindent\textbf{OE-15 Liquidación Interna de Cuentas.}
Implementar el módulo de liquidación interna para que los comensales puedan visualizar
el desglose de su consumo individual, los productos compartidos de la mesa y el monto
pendiente. El sistema permitirá registrar aportaciones realizadas por los comensales,
sin procesar transacciones bancarias reales ni integrarse con agregadores de pago
externos en la versión 1.0.

\textit{Evidencia: Visualización del desglose individual y compartido. Registro de
aportaciones internas. Validación de que el monto liquidado no exceda el total de la
cuenta. Cambio automático a estado \textit{Liquidada} cuando el total pendiente sea
cubierto.}

\noindent\textbf{OE-16 Consistencia Transaccional en Liquidaciones.}
Diseñar e implementar mecanismos de control para evitar inconsistencias durante el
registro de aportaciones simultáneas. El sistema deberá validar el estado actual de la
mesa, el monto pendiente y las aportaciones previamente registradas antes de aceptar
una nueva operación de liquidación. Para ello, se considerarán mecanismos de bloqueo
transaccional en PostgreSQL, como \texttt{SELECT FOR UPDATE}.

\textit{Evidencia: Intento de liquidación simultánea controlado. Rechazo de aportaciones
duplicadas o superiores al monto pendiente. Persistencia consistente del estado financiero
interno de la mesa.}

\noindent\textbf{Módulo 7 — Procesamiento Analítico con Apache Spark}

\noindent\textbf{OE-17 Extracción de Datos Transaccionales.}
Implementar procesos de extracción batch desde la base de datos transaccional hacia
la capa Bronze de la arquitectura analítica. Las extracciones deberán conservar los
datos relevantes de operación, tales como órdenes, productos, mesas, comensales,
estados de cocina, liquidaciones internas y estructura organizacional.

\textit{Evidencia: Extracciones batch ejecutadas correctamente. Datos crudos almacenados
en capa Bronze. Correspondencia validada entre registros transaccionales y archivos
generados.}

\noindent\textbf{OE-18 Transformación y Limpieza de Datos con Spark.}
Desarrollar jobs de Apache Spark para limpiar, normalizar y enriquecer los datos
extraídos en la capa Bronze, generando datasets Silver consistentes. Las transformaciones
deberán resolver formatos de fecha, relaciones entre entidades, cálculo de montos,
duraciones operativas y clasificación de eventos relevantes.

\textit{Evidencia: Jobs de Spark documentados. Datasets Silver generados. Validación de
reglas de transformación. Comparación entre datos crudos y datos procesados.}

\noindent\textbf{OE-19 Generación de Tablas Analíticas Gold.}
Construir tablas analíticas Gold a partir de los datasets Silver, enfocadas en indicadores
de negocio y operación restaurantera. Estas tablas deberán permitir consultar métricas
como ingresos por sucursal, productos más consumidos, actividad por mesa, tiempos de
preparación, velocidad de venta, comportamiento por franja horaria y desempeño operativo
por nivel organizacional.

\textit{Evidencia: Tablas Gold generadas. Métricas calculadas correctamente. Consultas
analíticas validadas contra datos transaccionales de origen.}

\noindent\textbf{OE-20 Estimaciones Exploratorias de Demanda.}
Implementar un componente analítico exploratorio para identificar patrones históricos
de consumo por producto, sucursal y franja horaria. Este componente no tendrá como
objetivo construir un sistema predictivo avanzado, sino generar estimaciones descriptivas
y tendencias iniciales que apoyen la interpretación del comportamiento de demanda.

\textit{Evidencia: Cálculo de tendencias por producto, sucursal y horario. Visualización
de patrones de consumo. Documentación de método utilizado y limitaciones del análisis.}

\noindent\textbf{Módulo 8 — Dashboard Analítico y Validación}

\noindent\textbf{OE-21 Dashboard Administrativo de Inteligencia de Negocio.}
Desarrollar un dashboard administrativo conectado a las tablas analíticas Gold para
visualizar indicadores relevantes de operación y negocio. El dashboard deberá permitir
filtrar información según el nivel organizacional del usuario, mostrando métricas
consolidadas para Admin Cadena, métricas regionales para Gerente Zona y métricas
locales para Gerente Sucursal.

\textit{Evidencia: Dashboard funcional con métricas provenientes de tablas Gold.
Filtros por nivel organizacional. Visualización de indicadores de ventas, productos,
mesas, tiempos operativos y demanda exploratoria.}

\noindent\textbf{OE-22 Validación de Métricas y Trazabilidad Analítica.}
Validar que las métricas mostradas en el dashboard correspondan con los datos
transaccionales registrados en el sistema. Para ello, se documentará la trazabilidad
entre eventos operativos, registros de base de datos, transformaciones Spark y tablas
analíticas finales.

\textit{Evidencia: Matriz de trazabilidad entre datos operativos y métricas analíticas.
Casos de prueba documentados. Comparación manual o automatizada entre registros
transaccionales y resultados agregados.}

\noindent\textbf{OE-23 Seguridad, Auditoría y Calidad del Sistema.}
Implementar mecanismos básicos de seguridad y calidad, incluyendo comunicación cifrada
mediante HTTPS, validación de accesos por rol, auditoría de operaciones críticas y
pruebas funcionales sobre los módulos principales. La auditoría deberá registrar eventos
administrativos y operativos relevantes para fortalecer la trazabilidad del sistema.

\textit{Evidencia: Aplicación operando sobre HTTPS. Pruebas funcionales documentadas.
Registros de auditoría generados. Validación de acceso por rol. Revisión de seguridad
sin vulnerabilidades críticas conocidas en los flujos principales.}

\newpage

\subsection{Justificación}

La relevancia de este proyecto radica en el diseño de una solución tecnológica que no
solo atiende la gestión operativa de mesas, órdenes y liquidaciones en restaurantes, sino
que además estructura dichos procesos como una fuente formal de datos transaccionales
para su posterior análisis. En este sentido, la plataforma propuesta funciona como un
sistema generador de eventos operativos, donde cada interacción del comensal, mesero,
cocina y administración queda registrada de manera consistente en una base de datos
centralizada.

El valor principal del trabajo se encuentra en la integración entre el modelo transaccional
y el modelo analítico. A partir de las órdenes registradas, los estados de preparación, las
sesiones de mesa, los consumos individuales y las liquidaciones internas, el sistema permite
construir un flujo de procesamiento batch mediante Apache Spark, organizado bajo una
arquitectura Medallion en capas Bronze, Silver y Gold. Esta arquitectura facilita la limpieza,
transformación y agregación de los datos para convertirlos en métricas útiles para la toma
de decisiones.

Desde la perspectiva de Ciencia de Datos, el proyecto permite aplicar conocimientos de
diseño de bases de datos, modelado transaccional, procesamiento distribuido, ingeniería de
datos, construcción de indicadores y visualización analítica. Las métricas generadas podrán
apoyar el análisis de ventas por sucursal, comportamiento de consumo, productos con mayor
demanda, tiempos de atención, desempeño de estaciones de cocina y patrones operativos
relevantes para la administración restaurantera.

A diferencia de una aplicación enfocada únicamente en la captura de pedidos, la propuesta
plantea un flujo completo donde la operación diaria se convierte en información estructurada
y explotable. De esta manera, el sistema no se limita a resolver una problemática inmediata
de división de cuentas, sino que habilita una base de conocimiento para análisis posteriores,
reportes administrativos y estimaciones exploratorias de demanda.

Asimismo, el desarrollo conserva un enfoque práctico y delimitado: la versión 1.0 no
procesará pagos bancarios reales ni integrará agregadores externos, sino que registrará el
estado interno de liquidación de las cuentas. Esta decisión reduce riesgos técnicos, legales
y de seguridad, permitiendo concentrar el trabajo terminal en el diseño de la base de datos,
la trazabilidad de la operación, el procesamiento analítico y la generación de inteligencia
de negocio.
informaci\'on, constituyendo una aplicaci\'on pr\'actica e integradora de los conocimientos adquiridos en la formaci\'on acad\'emica.
